% !TEX root = ../../../tree_notes.tex
% =====================================================================
%  15.2  The up/down pair: richer per-node states   sec:c20-updown
%  Exports: def:c20-parent-piece  thm:c20-parent-piece-correct
%           lst:c20-updown  fig:c20-parentpiece  ex:c20-updown-run
%           prob:c20-updown-solved  lst:c20-1156d-soln
%  Same two-DFS machinery as 15.1, but the state at a node is a PAIR that
%  the up pass must carry and the combine at each node is non-trivial.
%  Anchor problem CF 1156D (0-1-Tree): count ordered pairs (x,y) whose path
%  is 0*1* (never a 0-edge after a 1-edge). Answer = sum over transition
%  vertices of (z[v]*o[v] - 1), where z/o are per-label 0/1 component sizes
%  -- each computed by an INVERTIBLE two-DFS reroot. Worked O(n^2) vs O(n)
%  comparison. All overlays hand-computed on the running tree T.
% =====================================================================
\section{The Up/Down Pair: Richer States and the Parent-Side Piece}\label{sec:c20-updown}

\begin{previouslybox}
\begin{enumerate}[nosep, leftmargin=1.6em]
  \item In \secref{c20-twodfs} the state was a single number and $\Phi$ was a
        sum. What two facts let us skip a sibling loop? ($\Phi$ invertible, so
        "everything but $c$" $=$ subtract; and the edge-transfer had a closed
        form)
  \item \thmref{c20-edge-transfer} says $\mathit{up}[c]$ is built from
        $v$'s self, $\mathit{up}[v]$, and $c$'s siblings. If the state were a
        \emph{pair} $(a,b)$, what has to happen to \emph{both} components at
        every edge? (both are transferred; the pair is combined, not one)
  \item Ch.~13's independent-set in/out pair (\thmref{c10-mis-tree}) already
        carried two numbers per node. Was that a reroot? (no -- read at one
        fixed root; today we keep the pair but read it at \emph{every} node)
  \item If we naively re-rooted a per-node pair DP from all $n$ roots, what
        would the cost be, and what fixes it? ($\bigO{n^2}$; the two-DFS
        transfer, done once)
\end{enumerate}
\end{previouslybox}

\begin{topicintro}
Section 15.1 was the friendliest instance: one scalar, a sum, a closed-form
shift. Real problems hand the node a \emph{pair} of numbers -- an "answer from
below" and an "answer from above," or a "best within" and a "best crossing" --
and the interesting work is keeping both straight while the root slides. The
template does not change; only the meaning of $\mathit{down}[v]$,
$\mathit{up}[v]$ and their join grows. This section makes the parent-side piece
precise for a pair state (\defref{c20-parent-piece},
\thmref{c20-parent-piece-correct}), settles the running-cost question with a
worked $\bigO{n^2}$-versus-$\bigO{n}$ comparison, and instantiates everything
on CF~1156D: \emph{count the ordered pairs whose edge-label path is
$0^{\ast}1^{\ast}$}. There the pair is (size of $v$'s all-$0$ region, size of
$v$'s all-$1$ region), both are transferred by an invertible reroot, and $v$ is
scored as a \emph{transition vertex}.
\end{topicintro}

\subsection{The Parent-Side Piece, Precisely}

\begin{definition}{The pair state and its parent-side transfer}{c20-parent-piece}
Let the state at a node be a pair $\mathcal S(v) = (\mathit{down}[v],
\mathit{up}[v])$, where $\mathit{down}[v]$ is the fold over $\subtree(v)$ (the
Ch.~13 \thmref{c10-subtree-template} value) and $\mathit{up}[v]$ is the fold
over $V \setminus \subtree(v)$ assembled through $v$'s parent. The every-node
answer is a \emph{combine of the pair} at $v$,
$\mathit{dp}[v] = \Psi\bigl(\mathit{down}[v], \mathit{up}[v]\bigr)$ for some
node-local $\Psi$ (a sum, a product, a $\max$). When $\Phi$ is invertible the
parent-side piece obeys
\[
  \mathit{up}[c] = \mathit{shift}\!\Bigl(
      \underbrace{\mathit{up}[v] \;\oplus\; \text{self}(v)
      \;\oplus\; \bigl(\mathit{down}[v] \ominus \mathit{down}[c]\bigr)}_{\text{everything not in }\subtree(c)}
  \Bigr),
\]
where $\oplus$ is the combine, $\ominus$ its inverse (subtract / XOR-back), and
$\mathit{shift}$ the per-edge adjustment.
\end{definition}

The key pedagogical point is the term
$\mathit{down}[v] \ominus \mathit{down}[c]$: because the combine is a sum (or
XOR), the aggregate of $c$'s \emph{siblings} is not a loop over them -- it is
$v$'s own down-value with $c$'s contribution removed. That single subtraction
is what collapses \coderef{c20-twodfs}'s inner sibling loop from
$\bigO{\deg(v)}$ to $\bigO{1}$, and it is the difference between this problem
finishing in $\bigO{n}$ or in $\bigO{n^2}$ on a star.

\begin{theorem}{The pair reroot is correct and costs $\bigO{n}$ when $\Phi$ is invertible}{c20-parent-piece-correct}
For a state obeying \defref{c20-parent-piece} with invertible $\Phi$, the two
passes of \coderef{c20-updown} compute $\mathit{dp}[v] =
\Psi(\mathit{down}[v], \mathit{up}[v])$ for every node $v$ in $\bigO{n}$ time.
\end{theorem}

\begin{proofbox}[\thmref{c20-parent-piece-correct}: well-founded transfer, linear count]
\emph{Correctness.} By the same partition argument as
\thmref{c20-edge-transfer}: for each node $v$ the two-DFS computes the fold of
$\subtree(v)$ (down pass, by the induction hypothesis on Ch.~13's fold) and the
fold of the complement (up pass). Invertibility gives the siblings' fold exactly
as $\mathit{down}[v] \ominus \mathit{down}[c]$, so $\mathit{up}[c]$'s argument
is the true fold over $V\setminus\subtree(c)$, and $\mathit{shift}$ applies the
uniform per-edge re-indexing. The root seeds $\mathit{up}[r_0]=\Phi$'s identity,
correct because $V\setminus\subtree(r_0)=\varnothing$. Processing in pre-order
means $\mathit{up}[v]$ is finalised before it is used to form any
$\mathit{up}[c]$, so induction along the pre-order proves every $\mathit{dp}[v]$.
\emph{Cost.} Down pass: $\bigO{\deg(v)}$ per node, $\bigO{n}$ total. Up pass:
each node $v$ does $\bigO{1}$ work per child (one subtraction, one shift), so
$\bigO{\deg(v)}$ per node and $\bigO{\sum_v \deg(v)} = \bigO{n}$ total. No
sibling loop survives invertibility. \qed
\end{proofbox}

\subsection{The Cost of the Trick: All-Roots versus Two Passes}

\begin{center}
\small
\renewcommand{\arraystretch}{1.2}
\begin{tabular}{@{}p{0.42\linewidth} p{0.28\linewidth} p{0.24\linewidth}@{}}
\toprule
\textbf{Method} & \textbf{Time} & \textbf{Verdict at $n=2\cdot10^5$} \\
\midrule
Run Ch.~13's DFS once \emph{per root} & $\bigO{n \cdot n}=\bigO{n^2}$ & \tle \\
Two passes (down + up), invertible $\Phi$ & $\bigO{n}$ & \ac \\
Two passes, non-invertible $\Phi$, naive siblings & $\bigO{n + \sum\deg^2}$ (up to $\bigO{n^2}$ on a star) & \tle \\
Two passes, non-invertible, prefix/suffix (\secref{c20-prefixsuffix}) & $\bigO{n}$ & \ac \\
\bottomrule
\end{tabular}
\end{center}

\noindent The pair state does not change the \emph{shape} of the analysis --
it changes the constant (two numbers instead of one) and forces $\Psi$ at the
node. Invertibility is what keeps the up pass linear; the moment $\Psi$ or
$\Phi$ is a $\max$/product with no inverse, you owe yourself the prefix/suffix
argument of \secref{c20-prefixsuffix}. This is the seam between 15.2 and 15.3.

\subsection{The Algorithm: Pair Reroot with Subtract-the-Child}

\begin{lstlisting}[caption={The pair-state two-DFS: both components transfer, siblings recovered by inverting Phi (no O(deg) loop).}, label={lst:c20-updown}]
// Pair-state rerooting with INVERTIBLE combine (Phi = + here).
// down[v] = fold over subtree(v);  up[v] = fold over complement via parent.
// dp[v]   = Psi(down[v], up[v]) read at EVERY node.
vector<vector<int>> adj;
vector<Val> down_, up_, dp_;

Val combine(Val a, Val b);        // Phi  (associative, invertible)
Val inverse(Val a, Val b);        // a minus b under combine (e.g. subtraction)
Val shift(Val aggregate, int v, int c);
Val psi(Val d, Val u);            // node-local combine of the pair

void dfs_down(int v, int p) {                       // post-order
    down_[v] = self_value(v);
    for (int c : adj[v]) if (c != p) {
        dfs_down(c, v);
        down_[v] = combine(down_[v], down_[c]);
    }
}
void dfs_up(int v, int p) {                          // pre-order
    dp_[v] = psi(down_[v], up_[v]);                   // <-- answer for root v
    for (int c : adj[v]) if (c != p) {
        Val siblings = inverse(down_[v], down_[c]);   // down[v] minus down[c]
        Val allbut_c = combine(combine(up_[v], self_value(v)), siblings);
        up_[c] = shift(allbut_c, v, c);               // thm:c20-edge-transfer
        dfs_up(c, v);
    }
}
// seed: dfs_down(1,0); up_[1] = Phi_identity; dfs_up(1,0).
\end{lstlisting}

\begin{figure}[H]
\centering
\begin{tikzpicture}[diagram, scale=0.9, transform shape,
    boxnode/.style={rectangle, rounded corners=2pt, draw=#1, line width=0.8pt,
                    fill=#1!8, minimum width=15mm, minimum height=7mm,
                    font=\scriptsize\bfseries, text=#1!60!black}]
  \node[boxnode=successgreen] (d) at (0,0.9)  {$\mathit{down}[v]$};
  \node[boxnode=primaryblue]  (u) at (0,-0.9) {$\mathit{up}[v]$};
  \node[boxnode=difforange]   (p) at (4,0)    {$\mathit{dp}[v]=\Psi(\mathit{down},\mathit{up})$};
  \node[boxnode=warningred]   (s) at (2,1.9)  {siblings $=\mathit{down}[v]\ominus\mathit{down}[c]$};
  \draw[->, line width=0.8pt] (d) -- (p);
  \draw[->, line width=0.8pt] (u) -- (p);
  \draw[->, line width=0.8pt, dashed] (s) -- (u)
      node[midway, right, font=\scriptsize, text=warningred!70!black, anchor=west]
      {$\mathit{up}[c]$ uses self, $\mathit{up}[v]$, siblings};
  \node[annot, font=\scriptsize, text=successgreen!50!black, align=left]
    at (0,2.6) {down pass: post-order};
  \node[annot, font=\scriptsize, text=primaryblue!60!black, align=left]
    at (0,-2.1) {up pass: pre-order};
\end{tikzpicture}
\caption{The pair state at a node. The answer $\mathit{dp}[v]$ is a
node-local combine $\Psi$ of the two independent pieces. Only the up piece is
new relative to Chapter~13: to form $\mathit{up}[c]$ we need $v$'s self-value,
$v$'s own up-value, and $c$'s siblings -- and because $\Phi$ is invertible the
siblings are $\mathit{down}[v]\ominus\mathit{down}[c]$, a single subtraction,
not a loop.}\label{fig:c20-parentpiece}
\end{figure}

\subsection{Worked Example: the 0/1 Pair on $T$}

\begin{example}{Counting $0^{\ast}1^{\ast}$ ordered pairs on the running tree $T$}{c20-updown-run}
Reuse $T$ (root $1$). Give the nine edges the labels shown in the table; the
quantity at each node $v$ is the pair $(\mathit{down}0[v], \mathit{down}1[v])$
of \emph{all-$0$} and \emph{all-$1$} region sizes, and the node's contribution
to the answer is $z[v]\cdot o[v]-1$, where $z[v]$ (resp.\ $o[v]$) is $v$'s full
$0$-component (resp.\ $1$-component) size after the up pass. Every value is
hand-computed and reused unchanged by the trace box.
\end{example}

\begin{center}
\footnotesize
\renewcommand{\arraystretch}{1.2}
\begin{tabular}{l|cccccccccc}
\toprule
node $v$ & 1 & 2 & 3 & 4 & 5 & 6 & 7 & 8 & 9 & 10 \\
\midrule
labels on $(v,\text{parent})$ & -- & 0 & 1 & 0 & 1 & 0 & 1 & 1 & 0 & 0 \\
$z[v]$ ($0$-comp size) & 3 & 3 & 2 & 3 & 1 & 2 & 3 & 1 & 3 & 3 \\
$o[v]$ ($1$-comp size) & 3 & 3 & 3 & 1 & 3 & 1 & 3 & 3 & 1 & 1 \\
$z\!\cdot\!o-1$ & 8 & 8 & 5 & 2 & 2 & 1 & 8 & 2 & 2 & 2 \\
\bottomrule
\end{tabular}
\end{center}

\noindent Edge labelling used: $1\text{--}2{:}0$, $1\text{--}3{:}1$,
$2\text{--}4{:}0$, $2\text{--}5{:}1$, $3\text{--}6{:}0$, $3\text{--}7{:}1$,
$5\text{--}8{:}1$, $7\text{--}9{:}0$, $9\text{--}10{:}0$. Sanity check on the
$0$-edges $\{12,24,36,79,9\text{-}10\}$: they split $T$ into the components
$\{1,2,4\}$, $\{3,6\}$, $\{7,9,10\}$, $\{5\}$, $\{8\}$, which is exactly the
$z$ row above. The $1$-edges $\{13,25,37,58\}$ split $T$ into
$\{1,3,7\}$, $\{2,5,8\}$, and four singletons, giving the $o$ row. The node
contributions $z\cdot o-1$ then sum to $40$.

\begin{tracebox}[Recovering $z$ through the parent by invertible reroot, $0$-subproblem]
Root at $1$, $\mathit{down}0$ (0-only region inside subtree):
$\mathit{down}0[1]=3$ (nodes $1,2,4$), $\mathit{down}0[3]=2$ ($3,6$),
$\mathit{down}0[7]=3$ ($7,9,10$), $\mathit{down}0[5]=1$, $\mathit{down}0[8]=1$.
$z[1]=\mathit{down}0[1]=3$ (root has no parent side).\\
Transfer $z[c]=\mathit{down}0[c]+(\text{edge }0 \Rightarrow z[v]-\mathit{down}0[c])$:
$1\!\xrightarrow{0}\!2$: $z[2]=\mathit{down}0[2]+(z[1]-\mathit{down}0[2])=z[1]=3$.\\
$1\!\xrightarrow{1}\!3$: edge is 1, so no parent side: $z[3]=\mathit{down}0[3]=2$.\\
$2\!\xrightarrow{0}\!4$: $z[4]=z[2]=3$. \quad $2\!\xrightarrow{1}\!5$: $z[5]=\mathit{down}0[5]=1$.\\
$3\!\xrightarrow{0}\!6$: $z[6]=z[3]=2$. \quad $3\!\xrightarrow{1}\!7$: $z[7]=\mathit{down}0[7]=3$.\\
$5\!\xrightarrow{1}\!8$: $z[8]=\mathit{down}0[8]=1$. \quad $7\!\xrightarrow{0}\!9$: $z[9]=z[7]=3$.
\quad $9\!\xrightarrow{0}\!10$: $z[10]=z[9]=3$. \checkmark\ (matches $z$ row)
The $1$-subproblem is symmetric and yields the $o$ row. \checkmark
\end{tracebox}

\subsection{Variations and Twists}

\begin{itemize}[nosep]
  \item \textbf{Pair of counts (this problem):} both components are sizes; the
        node combine $\Psi$ is a product $z\!\cdot\!o$. \hfill\lv{4}
  \item \textbf{Best-within / best-crossing:} a max-plus pair (the diameter or
        longest-path state) -- $\Psi$ is a $\max$, still invertible for the
        \emph{transfer} if you keep the two top children. \hfill\lv{4}
  \item \textbf{Counting modulo a prime:} swap $-$ for modular inverse only if
        $\Phi$ is a product; for a sum keep plain subtraction and reduce.
        \hfill\lv{3}
  \item \textbf{Non-invertible $\Psi$:} the pair is fine, but the sibling
        recovery breaks -- that is \secref{c20-prefixsuffix}. \hfill\lv{4}
\end{itemize}

\subsection{Common Pitfalls}

\begin{warningbox}[Trap 1 -- combining the pair too early]
Form $\mathit{dp}[v]=\Psi(\mathit{down}[v],\mathit{up}[v])$ \emph{only after
both} passes are final for $v$. If you fold $\mathit{up}$ into $\mathit{down}$
during the down pass you lose the clean sibling inverse and the transfer reads
a double-counted quantity -- a symptom is the answer exceeding $\binom{n}{2}$
even before checking parity.
\end{warningbox}

\begin{warningbox}[Trap 2 -- a "0-edge after 1-edge" pair counted twice]
For $0^{\ast}1^{\ast}$ counting, the transition vertex is unique only when both
sides are non-empty; all-$0$ and all-$1$ paths must be assigned to a fixed end
(the formula $z\!\cdot\!o-1$ already does this by excluding the double-counted
self pair). If you instead count at \emph{both} endpoints of a monotone path
you over-count by exactly the pure-$0$ and pure-$1$ pairs.
\end{warningbox}

\begin{warningbox}[Trap 3 -- 32-bit sizes, forgotten long long product]
$z[v]\cdot o[v]$ is up to $n^2 \sim 4\cdot10^{10}$; the running sum is up to
$\Theta(n^2)$. Keep both $z,o$ as \lstinline{int} (sizes fit) but cast
\emph{before} multiplying: \lstinline{ans += 1LL*z[v]*o[v] - 1;}.
\end{warningbox}

\subsection{Problem Lab}

\solvedproblem{c20-updown-solved}{0-1-Tree}{Codeforces 1156D}
  {https://codeforces.com/problemset/problem/1156/D}{CF 2200 · L4}
  {pair-state rerooting, transition-vertex decomposition, invertible two-DFS
   component sizes, long long counting}

\begin{problemstatement}[0-1-Tree (statement and sample verified against the live problem)]
A tree of $n$ vertices ($2 \le n \le 2\cdot10^{5}$); every edge is labelled $0$
or $1$. An ordered pair $(x,y)$ with $x \ne y$ is \emph{valid} if, traversing
the simple path from $x$ to $y$, you never go through a $0$-edge \emph{after}
having gone through a $1$-edge -- i.e.\ the path's labels are $0^{\ast}1^{\ast}$.
Count the valid ordered pairs. \textbf{Example}: $n=7$, edges
$2\text{--}1{:}1$, $3\text{--}2{:}0$, $4\text{--}2{:}1$, $5\text{--}2{:}0$,
$6\text{--}7{:}1$, $7\text{--}2{:}1$; output \texttt{34}.
\end{problemstatement}

\begin{attemptgate}[Attempt it yourself first -- \lv{4}, give yourself 35 minutes]
Two structural facts. (a)~A $0^{\ast}1^{\ast}$ path has a unique
\emph{transition vertex} $m$ where the $0$-part ends and the $1$-part begins.
(b)~Everything on the $x$-side of $m$ is $0$-connected to $m$; everything on
the $y$-side is $1$-connected. So the count factors through $m$: how big is
$m$'s all-$0$ region, and its all-$1$ region?
\peek{Fix the transition vertex $m$. What set can $x$ be chosen from, and what
      set for $y$?}
\nudge{$x$ ranges over $m$'s $0$-edge-connected component, $y$ over $m$'s
      $1$-edge-connected component. Let $z[m], o[m]$ be those sizes. Which
      pairs are double-counted, and how do you exclude $x=m=y$?}
\fullhint{For each $m$, valid pairs that transition \emph{at} $m$ (x-side via
$0$-edges, y-side via $1$-edges, allowing either side to be just $m$ itself)
number $(z[m]-1)+(o[m]-1)+(z[m]-1)(o[m]-1) = z[m]\,o[m]-1$. Every $0^{\ast}1^{\ast}$
pair has a \emph{unique} such $m$ (for all-$0$ paths $m=y$; all-$1$ paths $m=x$),
so the answer is $\sum_m (z[m]\,o[m]-1)$. Compute $z$ as a per-node component
size by a two-DFS reroot on the $0$-edges only ($\Phi=+$, invertible):
$\mathit{down}0[v]=1+\sum_{0\text{-child }c}\mathit{down}0[c]$, then
$z[c]=\mathit{down}0[c]+(\text{edge }0?z[v]-\mathit{down}0[c]:0)$. Symmetric
$\mathit{down}1,o$ for the $1$-edges. Sum with \lstinline{long long}.}
\end{attemptgate}

\solstep{Understanding the Problem}
The condition "no $0$-edge after a $1$-edge" is a statement about the
\emph{shape} of the label sequence along the path: it is some (possibly empty)
run of $0$s followed by some (possibly empty) run of $1$s. Such a path is fully
described by the single vertex where the run of $0$s hands off to the run of
$1$s -- call it the transition vertex $m$. The portion before $m$ uses only
$0$-edges, so its far endpoint $x$ lies in $m$'s $0$-edge-connected component;
the portion after $m$ uses only $1$-edges, so $y$ lies in $m$'s
$1$-edge-connected component. This turns a global "for every pair, walk the
path" problem ($\bigO{n^2}$ paths) into a \emph{per-node} count -- exactly the
every-node aggregation the two-DFS template is built for, with a pair state
$(z[m],o[m])$.

\solstep{Key Observations}
\begin{observation}
A pair is valid iff its path's labels are $0^{\ast}1^{\ast}$, which holds iff
there is a vertex $m$ on the path splitting it into an all-$0$ prefix and an
all-$1$ suffix. $m$ is the transition vertex.
\end{observation}
\begin{observation}
For a fixed $m$: $x$ may be any vertex of $m$'s $0$-component and $y$ any
vertex of $m$'s $1$-component, and the resulting path is valid with transition
$m$. Counting $x\ne y$ gives $z[m]\,o[m]-1$ (the $-1$ drops the case $x=y=m$).
\end{observation}
\begin{observation}
Each valid pair has a \emph{unique} transition vertex ($0^{a}1^{b}$: if
$a,b\ge1$ the split is the unique $0/1$ boundary; if $b=0$ assign $m=y$; if
$a=0$ assign $m=x$), so summing $z[m]\,o[m]-1$ over all $m$ neither misses nor
double-counts a pair.
\end{observation}

\solstep{Think It Through With Me}
\begin{thinkbox}
\deadend{For every ordered pair, walk the path with binary lifting and check the
$0^{\ast}1^{\ast}$ condition.}
{$\bigO{n^2}$ pairs each costing $\bigO{\log n}$ -- hopeless at $n=2\cdot10^{5}$.
And LCA queries are Chapter~16 anyway; the real structure is the transition
vertex, which a path-scan never exploits.}
\deadend{Merge $0$- and $1$-edges into one DSU and count component products.}
{A single DSU conflates the two labels; you need the $0$-component \emph{and}
the $1$-component \emph{separately} at each node. Two DSUs (one per label) do
compute $z,o$ and are a valid alternative -- but the reroot below is the point
of this section and reuses \thmref{c20-edge-transfer} directly.}
\aha{Score each node as a transition vertex with the pair $(z[m],o[m])$; answer
$=\sum_m(z[m]\,o[m]-1)$. Get $z$ and $o$ by two invertible reroots (one on
$0$-edges, one on $1$-edges). That is \defref{c20-parent-piece} with $\Phi=+$.}
\end{thinkbox}

\solstep{Solution}
Root at $1$. For label $K\in\{0,1\}$: post-order
$\mathit{down}K[v]=1+\sum_{c:\,\text{edge}(v,c)=K}\mathit{down}K[c]$; pre-order
transfer $K[c]=\mathit{down}K[c]+(\text{edge}(v,c)=K?\;\mathit{full}K[v]-\mathit{down}K[c]:0)$
with $\mathit{full}K[1]=\mathit{down}K[1]$. Set $z=\mathit{full}0$,
$o=\mathit{full}1$. Answer $=\sum_{v} (z[v]\cdot o[v]-1)$ in
\lstinline{long long}.

\solstep{Proof of Correctness}
\begin{proofbox}[\coderef{c20-1156d-soln} counts exactly the valid ordered pairs]
\emph{Bijection.} Map a valid pair $(x,y)$ to its transition vertex $m$ as in
Observation~3; the map is well-defined and single-valued. Its inverse: given
$m$, a choice $x\in$ $0$-component of $m$, $y\in$ $1$-component of $m$, not both
$m$, produces the path $x\!\to\!m$ (all $0$) then $m\!\to\!y$ (all $1$), which
is $0^{\ast}1^{\ast}$ and whose transition (under the $b=0\Rightarrow m=y$,
$a=0\Rightarrow m=x$ conventions) is $m$. The two maps are mutual inverses, so
the number of valid pairs equals $\sum_m (z[m]o[m]-1)$. \emph{Component sizes.}
The reroot computes $z[v]=\card{\text{$0$-component of $v$}}$: down gives the
$0$-region within $\subtree(v)$; the transfer across a $0$-edge adds $v$'s side
by the inverse ($z[v]-\mathit{down}0[c]$ recovers exactly the $0$-region
outside $c$'s subtree), and across a $1$-edge contributes nothing (a $1$-edge is
not a $0$-path). By \thmref{c20-parent-piece-correct} every $z[v]$ is correct;
symmetrically for $o$. \qed
\end{proofbox}

\solstep{Implementation}
\begin{lstlisting}[caption={Codeforces 1156D -- two invertible pair reroots (0- and 1-edges), then sum z[v]*o[v] - 1 over every transition vertex}, label={lst:c20-1156d-soln}]
#include <bits/stdc++.h>
using namespace std;
typedef long long ll;
int n;
struct Edge { int to, c; };
vector<vector<Edge>> adj;
vector<int> d0, d1;    // down: all-0 / all-1 region within subtree
vector<int> z, o;      // full 0-component / 1-component sizes per node

void dfs_down(int v, int p, vector<int>& d, int K) {   // post-order
    d[v] = 1;
    for (const Edge& e : adj[v]) if (e.to != p) {
        dfs_down(e.to, v, d, K);
        if (e.c == K) d[v] += d[e.to];                 // fold same-label child
    }
}
void dfs_up(int v, int p, vector<int>& d, vector<int>& full, int K) {
    for (const Edge& e : adj[v]) if (e.to != p) {       // pre-order
        int add = (e.c == K) ? full[v] - d[e.to] : 0;   // parent side via K-edge
        full[e.to] = d[e.to] + add;
        dfs_up(e.to, v, d, full, K);
    }
}
int main() {
    scanf("%d", &n);
    adj.assign(n + 1, {});
    for (int i = 0; i < n - 1; i++) {
        int x, y, c; scanf("%d %d %d", &x, &y, &c);
        adj[x].push_back({y, c}); adj[y].push_back({x, c});
    }
    d0.assign(n+1,0); d1.assign(n+1,0); z.assign(n+1,0); o.assign(n+1,0);
    dfs_down(1, 0, d0, 0); dfs_down(1, 0, d1, 1);
    z[1] = d0[1]; o[1] = d1[1];
    dfs_up(1, 0, d0, z, 0); dfs_up(1, 0, d1, o, 1);
    ll ans = 0;
    for (int v = 1; v <= n; v++) ans += 1LL * z[v] * o[v] - 1;
    printf("%lld\n", ans);
    return 0;
}
\end{lstlisting}

\solstep{Trace}
\begin{tracebox}[Official Example: $n=7$, edges $2\text{--}1{:}1,3\text{--}2{:}0,4\text{--}2{:}1,5\text{--}2{:}0,6\text{--}7{:}1,7\text{--}2{:}1$]
Root at $1$. $0$-edges $=\{3\text{--}2,5\text{--}2\}$, so the $0$-component
$\{2,3,5\}$ has size $3$; every other node's $0$-component is a singleton.
$z=[1,3,3,1,3,1,1]$ for $v=1..7$. $1$-edges $=\{2\text{--}1,4\text{--}2,
6\text{--}7,7\text{--}2\}$ connect $\{1,2,4,6,7\}$ (size $5$); nodes $3,5$ have
no $1$-edge. $o=[5,5,1,5,1,5,5]$. Contributions $z\!\cdot\!o-1$:
$1{\cdot}5{-}1{=}4$, $3{\cdot}5{-}1{=}14$, $3{\cdot}1{-}1{=}2$,
$1{\cdot}5{-}1{=}4$, $3{\cdot}1{-}1{=}2$, $1{\cdot}5{-}1{=}4$,
$1{\cdot}5{-}1{=}4$. Sum $=4{+}14{+}2{+}4{+}2{+}4{+}4=\boxed{34}$ \checkmark
matches CF's example.
\end{tracebox}

\solstep{Complexity}
\begin{complexitybox}
Time $\bigO{n}$: two down passes and two up passes, each $\bigO{\sum\deg}
=\bigO{n}$, since $\Phi=+$ is invertible and each transfer is one subtraction --
no sibling loop (\thmref{c20-parent-piece-correct}). Memory $\bigO{n}$ (four
arrays plus adjacency). $z[v]o[v]-1$ summed over $n$ nodes reaches
$\bigO{n^2}\sim 4\cdot10^{10}$, hence \lstinline{long long} for the accumulator.
Recursion depth is $h$; at $n\le2\cdot10^{5}$ with a chain legal, run the four
passes iteratively (\coderef{c09-iter}) or the run \re's. Verdict: \ac; the
archetypal \emph{pair-state} reroot -- two independent invertible up/down
passes combined by a node-local product.
\end{complexitybox}

\solstep{Pitfalls and Variants}
\begin{warningbox}
Ways to lose CF~1156D: (i)~\lstinline{int} accumulator -- overflow past $n^2$;
(ii)~letting a $0$-component "leak" across a $1$-edge in the transfer (the
\lstinline{e.c == K} guard is exactly the $0^{\ast}1^{\ast}$ boundary);
(iii)~counting the transition at both ends and reading $\sum(z\,o-1)$ as if it
were a plain component-product -- the $-1$ and the unique-$m$ argument are load-
bearing; (iv)~forgetting $x\ne y$. Variant: if the condition becomes "at most
one $1$-edge," the same pair state works with $o$-components capped at depth~1;
if it becomes "$0^{\ast}1^{\ast}0^{\ast}$" the combine is no longer a single
product and \secref{c20-prefixsuffix} is required.
\end{warningbox}

\subsection{Practice Problems}

\begin{stuckbox}
\begin{itemize}[nosep, leftmargin=1.4em]
  \item Did you identify the \emph{transition vertex} so the pair $(z,o)$ is a
        per-node quantity, before writing any reroot?
  \item Are both $0$- and $1$-regions computed by \emph{separate} invertible
        reroots, with the same-label guard in the transfer?
  \item Is the final sum in \lstinline{long long}, subtracting $1$ per node for
        the excluded $x=y$ pair?
\end{itemize}
\end{stuckbox}

\begin{practiceset}
  \practice{Three Paths on a Tree}{Codeforces 1294F}
    {https://codeforces.com/problemset/problem/1294/F}{CF 2000 · L4}
    {Choose three distinct vertices $a,b,c$ maximising the number of edges in
    the union of the three pairwise paths. The union is the minimal subtree on
    $\{a,b,c\}$, shaped like a three-armed star meeting at one centre; its size
    is $\tfrac12(\dist(a,b)+\dist(b,c)+\dist(a,c))$. Take a diameter's endpoints
    as two of $a,b,c$ and the third as the vertex farthest from the diameter
    path -- a rerooting of "deepest arm off each node" ($\mathit{down}$ from the
    children, $\mathit{up}$ through the parent) on top of the two-BFS diameter
    of Ch.~12. Samples (verified): the $8$-node sample outputs $5$ and a triple
    such as \texttt{1 8 6}.
    }
\end{practiceset}

\paragraph{Back-cites (URLs live elsewhere in the book).}
\begin{itemize}[nosep,leftmargin=1.4em]
  \item \textbf{CF~543D ``Road Improvement''} and \textbf{CF~1092F ``Tree with
        Maximum Cost''} -- both are pair/count reroots of this shape; their URLs
        are Chapter~12's / the mixed set's, so cite by prose, not by link here.
\end{itemize}

\begin{recallbox}[10-second refresher -- 15.2 The up/down pair]
When the state is a pair $(\mathit{down}[v],\mathit{up}[v])$ and $\Phi$ is
invertible, the up pass recovers a child's siblings as
$\mathit{down}[v]\ominus\mathit{down}[c]$ -- one subtraction, not a loop -- so
the reroot is $\bigO{n}$ (\thmref{c20-parent-piece-correct}). CF~1156D
(\probref{c20-updown-solved}) scores each node as a $0^{\ast}1^{\ast}$
transition vertex: pair $(z[v],o[v])$ of per-label component sizes, answer
$\sum_v(z[v]o[v]-1)$. On $T$ with the opener's labelling the total is $40$; on
the official $n=7$ example it is $34$.
\selftest{
  \item Which single term in \defref{c20-parent-piece} turns the naive
        $\bigO{\deg(v)}$ sibling aggregate into $\bigO{1}$?
  \item Why is the transition vertex of a $0^{\ast}1^{\ast}$ path unique, and
        how are the all-$0$ / all-$1$ paths assigned?
  \item State the transfer $z[c]$ for a $0$-edge versus a $1$-edge.
  \item What does the $-1$ in $z[v]o[v]-1$ exclude?
  \item If the combine became a $\max$, would invertibility still hold -- and
        what section do you owe then?
}
\end{recallbox}
